La arquitectura del sistema sigue el modelo clásico de compilación y ejecución de programas.

\section{Análisis Léxico y Sintáctico}
El compilador utiliza \textit{Flex} y \textit{Bison} para la generación automática de los analizadores:
\begin{itemize}
    \item \textbf{Flex (src/lexer.l):} Define los patrones léxicos del lenguaje (tokens). Convierte el flujo de caracteres de entrada en una secuencia de tokens.
    \item \textbf{Bison (src/parser.y):} Define la gramática libre de contexto del lenguaje. Construye un Árbol de Sintaxis Abstracta (AST) que representa la estructura jerárquica del programa.
\end{itemize}

\section{Análisis Semántico y Generación de Código}
Una vez construido el AST, el compilador realiza:
\begin{enumerate}
    \item \textbf{Tabla de Símbolos (src/TablaSimbolos.cc):} Gestiona el ámbito y el tipo de los identificadores declarados.
    \item \textbf{Tabla de Tipos (src/TablaTipos.cc):} Verifica la consistencia de los tipos en las expresiones (verificación de tipos).
    \item \textbf{Generador de Código:} Recorre el AST para emitir las instrucciones de bajo nivel compatibles con la VM \textit{m2r}.
\end{enumerate}

\section{Máquina Virtual m2r}
El intérprete (\texttt{m2r.c}) emula una arquitectura basada en:
\begin{itemize}
    \item \textbf{Memoria:} Un array de estructuras \texttt{MEMO} que almacenan enteros o reales.
    \item \textbf{Registros:} Acumulador y Registro Base.
    \item \textbf{Ciclo de Instrucción:} \textit{Fetch-Decode-Execute}. Lee la instrucción en el PC (Program Counter), decodifica la operación y sus operandos, y ejecuta la acción correspondiente (aritmética, saltos, E/S).
\end{itemize}
